\documentclass[11pt]{article}
\usepackage[T1]{fontenc}
\usepackage{lmodern}
\usepackage[margin=1.05in]{geometry}
\usepackage{amsmath,amssymb,amsthm,mathtools}
\usepackage{microtype,booktabs,array,longtable}
\usepackage[colorlinks=true,linkcolor=blue!45!black,citecolor=blue!45!black,urlcolor=blue!45!black]{hyperref}
\usepackage{xcolor}
\usepackage[nameinlink,capitalise,noabbrev]{cleveref}
\usepackage{listings}
\lstset{basicstyle=\ttfamily\small,columns=fullflexible,keepspaces=true,breaklines=true,
  literate={≠}{{$\neq$}}1 {→}{{$\to$}}1 {∀}{{$\forall$}}1 {∃}{{$\exists$}}1 {¬}{{$\neg$}}1
           {∈}{{$\in$}}1 {∧}{{$\wedge$}}1 {∨}{{$\vee$}}1 {⟨}{{$\langle$}}1 {⟩}{{$\rangle$}}1}
\numberwithin{equation}{section}
\newtheorem{theorem}{Theorem}[section]
\newtheorem{proposition}[theorem]{Proposition}
\newtheorem{lemma}[theorem]{Lemma}
\newtheorem{corollary}[theorem]{Corollary}
\theoremstyle{definition}
\newtheorem{definition}[theorem]{Definition}
\theoremstyle{remark}
\newtheorem{remark}[theorem]{Remark}
\newcommand{\Z}{\mathbb Z}
\newcommand{\F}{\mathbb F}
\newcommand{\GF}{\mathrm{GF}}
\DeclareMathOperator{\supp}{supp}
\newcommand{\sds}[4]{\mathrm{SDS}(#1,#2,#3,[#4])}
\newcommand{\certrow}[5]{#1 & #2 & #3 & {\scriptsize #4} & {\scriptsize #5}\\[2pt]}
\title{Settling small open cases of signed difference sets}
\author{Claude (AI model by Anthropic)\\ \small run and submitted by Aleksandr Bersenev}
\date{October 8, 2026\\[3pt]\footnotesize Produced autonomously by Claude (Opus 5.5) agents in Claude Code; run and submitted by Aleksandr Bersenev.\\ \footnotesize Not an official Anthropic publication. Not peer reviewed.}
\hypersetup{pdftitle={Settling small open cases of signed difference sets},pdfauthor={Claude (AI model by Anthropic); submitted by Aleksandr Bersenev}}
\begin{document}
\maketitle
\begin{abstract}
A $(v,k,\lambda)$ signed difference set in a finite abelian group $G$ of order $v$ is a group-ring
element $D=\sum_g s_g g$ with $s_g\in\{-1,0,1\}$, exactly $k$ nonzero coefficients, and
$DD^{(-1)}=(k-\lambda)1_G+\lambda G$. Gordon's La Jolla signed difference set repository lists
67{,}823 group-specific parameter cases as open; after the nonexistence criteria of He and Wu
(2026), 22{,}462 remain. We prove that there is no $(9,8,1)$ signed difference set in
$\Z_3\times\Z_3$, the smallest open case; the proof is short, and a Lean~4 kernel-checked
certificate is given. We exhibit explicit signed difference sets for 45 further open cases, all with
$v\le 64$, each verified by the Lean kernel. To our knowledge, as of October 8, 2026, 26 of these are
new. These include an eighth-order cyclotomic $(25,12,1)$ signed difference set in $\GF(25)$ and
examples in $\Z_3\times\Z_9$, in six groups of order 32, and in seven groups of order 64.
Of eleven further cases of order 64, with parameters $(64,22,6)$ and $(64,46,30)$, two follow from
the partial-difference-set argument with partial difference sets that exist in the group, while
for the other nine the required partial difference sets do not exist in the group (exhaustive
search), so that these nine are new in the same sense. Eight are corrections to the repository,
since they follow from published constructions. Finally, an exhaustive computer search shows, without formal verification, that the
remaining 16 open cases with $v\le 28$ admit no signed difference set. Together with the results
above, this decides all 23 open cases with $v\le 28$.
\end{abstract}
\tableofcontents

%======================================================================
\section{Introduction}

Signed difference sets were introduced by Gordon~\cite{Gordon2023} as a common generalization of
difference sets (all coefficients $+1$) and group-invariant weighing matrices ($\lambda=0$).
Equivalently, in the cyclic case they are ternary periodic sequences with two-level autocorrelation.
Gordon gave several constructions, ran multiplier-based computer searches, and set up the
\emph{La Jolla Signed Difference Set Repository}~\cite{GordonSDSrepo}. For every admissible
triple $(v,k,\lambda)$ and every isomorphism type of abelian group of order $v<500$, the repository
records whether a signed difference set is known, known not to exist, or open. He, Chen and
Ge~\cite{HeChenGe2024} gave further constructions from Paley-type partial difference sets,
product constructions in $3$-groups, and fourth-order cyclotomy. Version~1.2 of the repository
(April 24, 2026) incorporates their Paley-type examples (entries marked ``PDS construction of He,
Chen and Ge''), but not their $(3^{2m+1},3^{2m}+1,1)$ family: $\sds{243}{82}{1}{3,3,3,3,3}$ is
still listed as open. Versions~1.0 (December 2022) and 1.1 (March 2023) of the repository covered
essentially only cyclic groups (1951 and 6719 entries, of which 1 and 28 are noncyclic); the
group-specific entries for noncyclic abelian groups, and with them every case treated in this
paper, were added in version~1.2, which lists $67{,}823$ group-specific cases as open. He and Wu~\cite{HeWu2026} proved four nonexistence criteria. Applied cumulatively, the
criteria exclude $45{,}361$ of the open cases and leave $22{,}462$. He and Wu also gave three new
existence results. We reran their verification script~\cite{HeScript}, and all published counts
were reproduced. The smallest remaining open cases are
\[
 \sds{9}{8}{1}{3,3},\ \sds{20}{11}{2}{2,10},\ \sds{20}{17}{8}{2,10},\ \sds{24}{18}{2}{2,12},\ \dots
\]
Here $[a_1,\dots,a_r]$ denotes $\Z_{a_1}\times\dots\times\Z_{a_r}$.

\paragraph{Results.} The results fall into five groups.
\begin{enumerate}
\item[(d)] \textbf{Nonexistence, Lean-certified.} There is no $(9,8,1)$ signed difference set in
$\Z_3\times\Z_3$ (\cref{thm:981}). The proof is a short intersection-number argument. The
statement is also checked by the Lean kernel by exhaustive enumeration, without symmetry
reductions.
\item[(a)] \textbf{New existence results.} \Cref{thm:new} gives explicit signed difference sets for 26
open cases. To our knowledge, as of October 8, 2026, none of them is implied by a published
construction.
\item[(b)] \textbf{Existence, order 64, with a partial-difference-set analysis.} For 11 open cases of
order 64 we give certified examples. Such cases could in principle follow from the
partial-difference-set argument of \cref{lem:pds} applied to Latin-square-type or
negative-Latin-square-type partial difference sets. An exhaustive search shows that the required
partial difference sets exist in only two of the eleven groups, so nine of these cases cannot
arise in this way (\cref{prop:maybe}).
\item[(c)] \textbf{Database corrections.} Eight cases listed as open follow from published
constructions or from routine variants of them (\cref{prop:known}). We do not claim them as new.
One of them is He and Wu's own $(64,22,6)$ in $\Z_2^6$.
\item[(e)] \textbf{Computer-only nonexistence.} An exhaustive search, which is \emph{not} formally
verified, shows that the remaining 16 open cases with $v\le 28$ have no signed difference set
(\cref{sec:nonex}). Two of them also have a short parity proof (\cref{prop:parity}).
\end{enumerate}
Consequently, all 23 open cases with $v\le 28$ are now decided: 6 exist and 17 do not
(\cref{tab:small}). All 45 existence certificates are listed in \cref{app:certs} and in the
file \texttt{data/certificates.json}. The Lean development is described in \cref{sec:lean}.
It uses Lean~4.34.1 with core only and no Mathlib, and its proofs depend only on the standard
axioms \texttt{propext} and \texttt{Quot.sound}.

\paragraph{How the examples were found.} The search used simulated annealing over all 710 open
cases with $v\le100$ (\cref{sec:methods}). It succeeded only for $v\le 64$. For the $(25,12,1)$
example we then found an algebraic description, as a signed union of eighth-order cyclotomic
classes of $\GF(25)$.

%======================================================================
\section{Preliminaries}\label{sec:prelim}

Let $G$ be a finite abelian group of order $v$, written additively, and let $\Z[G]$ be its group
ring. For $A=\sum_g a_g g$ put $A^{(-1)}=\sum_g a_g(-g)$. We identify a subset $S\subseteq G$ with
$\sum_{g\in S}g$.

\begin{definition}[\cite{Gordon2023}]
Let $D=\sum_{g\in G}s_g g$ with $s_g\in\{-1,0,1\}$. Put $P=\{g:s_g=1\}$, $N=\{g:s_g=-1\}$ and
$k=|P|+|N|$. Then $D=P-N$ is a \emph{$(v,k,\lambda)$ signed difference set} (SDS) in $G$ if
\begin{equation}\label{eq:sds}
DD^{(-1)}=(k-\lambda)\cdot 0_G+\lambda G .
\end{equation}
\end{definition}

The coefficient of $y$ in $DD^{(-1)}$ is the periodic autocorrelation
$C_D(y)=\sum_{g\in G}s_g s_{g+y}$. The coefficient at $y=0$ is $\sum_g s_g^2=k$ automatically.
So \eqref{eq:sds} is equivalent to
\begin{equation}\label{eq:autocorr}
C_D(y)=\sum_{g\in G}s_gs_{g+y}=\lambda\qquad\text{for all }y\in G\setminus\{0\}.
\end{equation}

\begin{lemma}[basic facts]\label{lem:basic}
Let $D=P-N$ be a $(v,k,\lambda)$ SDS in $G$.
\begin{enumerate}
\item $s^2=k+\lambda(v-1)$, where $s=|P|-|N|$~\cite{Gordon2023}.
\item $-D$, every translate $g+D$, and every image $\alpha(D)$ under an automorphism $\alpha$ of $G$
are again $(v,k,\lambda)$ SDSs.
\item (Support parity.) Let $S=P\cup N$ and $c_y=|S\cap(S-y)|$. Then $C_D(y)$ is a sum of $c_y$ terms
$\pm1$. Hence $|\lambda|\le c_y$ and $c_y\equiv\lambda\pmod 2$ for all $y\ne0$.
\end{enumerate}
\end{lemma}
\begin{proof}
(1) Apply the principal character to \eqref{eq:sds}. (2) The products $(-D)(-D)^{(-1)}$,
$(g+D)(g+D)^{(-1)}$ and $\alpha(D)\alpha(D)^{(-1)}=\alpha(DD^{(-1)})$ are all equal to
$(k-\lambda)\cdot0_G+\lambda G$. (3) The nonzero terms of $C_D(y)$ are those with $g\in S$ and
$g+y\in S$.
\end{proof}

We use the quotient method in the form of He and Wu~\cite[Lemma~4.1]{HeWu2026}. Let $H\le G$ with
$|H|=h$. For each coset $x$ of $H$ put $A_x=|P\cap x|-|N\cap x|$ and $K_x=|P\cap x|+|N\cap x|$.
Then
\begin{equation}\label{eq:quot}
A_x\equiv K_x\pmod 2,\qquad \sum_{x\in G/H}A_x^2=(k-\lambda)+\lambda h .
\end{equation}
This follows by mapping \eqref{eq:sds} to $\Z[G/H]$ and reading off the coefficient of the
identity.

\begin{lemma}[PDS argument; cf.\ the proof of {\cite[Thm.~5.3]{HeWu2026}}]\label{lem:pds}
Let $P\subseteq G\setminus\{0\}$ be a regular $(v,k_P,\lambda_P,\mu_P)$ partial difference set,
that is, $P=-P$ and $PP^{(-1)}=k_P\cdot 0_G+\lambda_PP+\mu_P(G-P-0_G)$. If $\lambda_P-\mu_P=2$,
then $D=P-\{0\}$ is a $(v,k_P+1,\mu_P)$ SDS.
\end{lemma}
\begin{proof}
Using $P^{(-1)}=P$ we get $DD^{(-1)}=P^2-2P+0_G=(k_P-\mu_P+1)\cdot0_G+(\lambda_P-\mu_P-2)P+\mu_P G$.
The middle term vanishes, so $D$ is an SDS with $\lambda=\mu_P$ and $k-\lambda=k_P-\mu_P+1$.
\end{proof}

%======================================================================
\section{There is no \texorpdfstring{$(9,8,1)$}{(9,8,1)} SDS in \texorpdfstring{$\Z_3\times\Z_3$}{Z3 x Z3}}\label{sec:981}

\begin{theorem}\label{thm:981}
There is no $(9,8,1)$ signed difference set in $\Z_3\times\Z_3$.
\end{theorem}
\begin{proof}
Suppose $D=P-N$ is one. By \cref{lem:basic}(1), $s^2=8+8=16$, so $|P|-|N|=\pm4$. Replacing
$D$ by $-D$ if necessary, we may assume $|P|=6$ and $|N|=2$. Exactly one element has coefficient
$0$. After a translation (\cref{lem:basic}(2)) we may assume this element is $0$.

Let $H$ be any of the four subgroups of order~$3$, and apply \eqref{eq:quot} with $h=3$. This
gives $A_H^2+A_{x}^2+A_{x'}^2=7+3=10$, where $x,x'$ are the other two cosets. The coset $H$
contains the zero coefficient, so $K_H=2$, while $K_x=K_{x'}=3$. Hence $A_H$ is even and
$A_x,A_{x'}$ are odd. If $A_H=\pm2$ then $A_x^2+A_{x'}^2=6$, which is impossible because odd
squares lie in $\{1,9,25,\dots\}$. Therefore $A_H=0$. The two nonzero elements of $H$ thus carry
opposite signs, and exactly one of them lies in $N$. The four subgroups of order~$3$ partition
$G\setminus\{0\}$, so $|N|=4$, which contradicts $|N|=2$.
\end{proof}

\begin{remark}
The case is not covered by the criteria of He and Wu. In their near-full-support test
\cite[Thm.~4.4]{HeWu2026}, with $h=q=3$, one gets $v-1+\lambda(h-1)=10\equiv 2\equiv q-1\pmod 8$,
which is allowed. The small-defect test \cite[Thm.~4.7]{HeWu2026} gives the same congruence. The
self-conjugacy test does not apply because $n=k-\lambda=7$ and $7\equiv1\pmod 3$. The
Lean proof in \cref{sec:lean} uses no symmetry reduction: it enumerates all $3^9$ coefficient
vectors.
\end{remark}

%======================================================================
\section{Existence results}\label{sec:exist}

The certificates in this section are explicit pairs $(P,N)$. Each is listed in \cref{app:certs},
in the file \texttt{data/\allowbreak certificates.json}, and verbatim in \texttt{lean/\allowbreak Main.lean}. Checking
\eqref{eq:autocorr} is a finite computation: $v\le 64$ products for each of the $v-1$ nonzero shifts.
For every certificate it was carried out in three ways: (i) by the Lean kernel (\cref{sec:lean});
(ii) by an independent Python script (\texttt{data/\allowbreak verify\_\allowbreak certificates.py}), which works with
coordinate tuples rather than the Lean encoding; (iii) during the search itself.

\subsection{New cases}

\begin{theorem}\label{thm:new}
In each of the following 26 groups there is an SDS with the stated parameters. The certificates
are in \cref{tab:certs} (tier~a).
\begin{center}
\begin{tabular}{ll}
\toprule
$(v,k,\lambda)$ & groups\\ \midrule
$(25,12,1)$ & $\Z_5^2$\\
$(27,10,1)$ & $\Z_3\times\Z_9$\\
$(27,14,5)$ & $\Z_3\times\Z_9$\\
$(32,20,4)$ & $\Z_2^5,\ \Z_2^3\times\Z_4,\ \Z_2^2\times\Z_8,\ \Z_2\times\Z_4^2,\ \Z_2\times\Z_{16},\ \Z_4\times\Z_8$\\
$(32,28,12)$ & $\Z_2\times\Z_4^2,\ \Z_4\times\Z_8$\\
$(36,11,2)$ & $\Z_3\times\Z_{12}$\\
$(48,24,8)$ & $\Z_2^3\times\Z_6,\ \Z_2^2\times\Z_{12},\ \Z_4\times\Z_{12}$\\
$(48,32,16)$ & $\Z_2^3\times\Z_6,\ \Z_2^2\times\Z_{12},\ \Z_4\times\Z_{12}$\\
$(54,46,10)$ & $\Z_3^2\times\Z_6$\\
$(64,18,2)$ & $\Z_2^6,\ \Z_2^4\times\Z_4,\ \Z_2^3\times\Z_8,\ \Z_2^2\times\Z_4^2,\ \Z_2\times\Z_4\times\Z_8,\ \Z_4^3,\ \Z_8^2$\\
\bottomrule
\end{tabular}
\end{center}
All 26 cases are listed as open in the repository~\cite{GordonSDSrepo}, and none is excluded by
\cite{HeWu2026}. To our knowledge, as of October 8, 2026, none of them follows from the
constructions in \cite{Gordon2023,HeChenGe2024,HeWu2026} or from \cref{lem:pds}.
\end{theorem}
\begin{proof}
For each case, the certificate in \cref{app:certs} satisfies \eqref{eq:autocorr}, which is a
finite computation, verified as described above.
\end{proof}

\begin{proposition}[an eighth-order cyclotomic SDS]\label{prop:cyc25}
Let $\GF(25)=\F_5[x]/(x^2+2)$, let $w=1+x$ (a primitive element), and let
$C_i=w^i\langle w^8\rangle$ for $0\le i<8$, the eighth-order cyclotomic classes ($|C_i|=3$). Then
\[
D=C_0-C_2+C_3+C_4
\]
is a $(25,12,1)$ SDS in $(\GF(25),+)\cong\Z_5^2$.
\end{proposition}
\begin{proof}
In the basis $\{1,x\}$ this $D$ is exactly the certificate of the Lean theorem
\texttt{sds\_25\_12\_1\_G5x5}, so \eqref{eq:autocorr} holds by finite computation. Since
$D$ is invariant under multiplication by $\langle w^8\rangle$, $C_D$ is constant on each class
$C_i$. It therefore suffices to check the eight values $C_D(w^i)$, $0\le i<8$, which all equal~1.
\end{proof}

\begin{remark}
He, Chen and Ge classified fourth-order cyclotomic SDSs~\cite{HeChenGe2024}. Gordon's cyclic
sporadic $(89,33,1)$~\cite{Gordon2023} is also a signed union of eighth-order cyclotomic classes
of $\GF(89)$. A search over all signed unions of eighth-order classes, possibly together with
$\pm\{0\}$, in $\GF(q)$ for $q\in\{9,17,25,41,49,73,89,97,113,121\}$ finds, apart from the trivial
parameters $(q,q-1,q-2)$ and the Paley-type $(q,q-1,-1)$ sets, only $(73,9,1)$ and its complement
$(73,64,56)$ (difference sets), $(73,64,0)$ (a circulant weighing matrix), $(73,36,4)$ and
$(89,33,1)$ (both in Gordon's table), and the new $(25,12,1)$. The local-search example for $(25,12,1)$ has
affine stabilizer of order~3 in $\mathrm{AGL}(2,5)$, and this is what led to the cyclotomic
description.
\end{remark}

\begin{remark}[structure of other examples]
(1) An exhaustive search shows that $\Z_3\times\Z_9$ contains no $(27,13,6)$ difference set.
So the route via Gordon's Prop.~2.3 (a skew Hadamard difference set plus $N=\{0\}$) is not
available for $(27,14,5)$ in $\Z_3\times\Z_9$, unlike in $\Z_3^3$ (see \cref{prop:known}).
(2) In both $(32,28,12)$ examples the zero set is a coset of a subgroup of order~4, namely of
$\{0,(0,0,2),(0,2,0),(0,2,2)\}$ in $\Z_2\times\Z_4^2$ and of $\{0,(0,4),(2,0),(2,4)\}$ in
$\Z_4\times\Z_8$; so after a translation the support is the complement of a subgroup of order~4.
(3) In $\Z_2^6$, the support of our $(64,18,2)$ example, translated by $(0,1,0,0,1,0)$ so that
it avoids $0$, is a $(64,18,2,6)$ partial difference set; it carries $15$ plus signs and $3$ minus
signs.
(4) The families $(32,20,4)$, $(48,24,8)$, $(48,32,16)$, $(64,18,2)$, $(64,22,6)$ and
$(64,46,30)$ all have $n=k-\lambda=16$. We do not know a uniform construction for them.
\end{remark}

\subsection{Cases of order 64 and the partial-difference-set route}

\begin{proposition}\label{prop:maybe}
SDSs exist in the following 11 cases (\cref{tab:certs}, tier~b):
$(64,22,6)$ in $\Z_2^4\times\Z_4$, $\Z_2^3\times\Z_8$, $\Z_2\times\Z_4\times\Z_8$, $\Z_4^3$,
$\Z_4\times\Z_{16}$; and $(64,46,30)$ in $\Z_2^4\times\Z_4$, $\Z_2^3\times\Z_8$,
$\Z_2^2\times\Z_4^2$, $\Z_2\times\Z_4\times\Z_8$, $\Z_4^3$, $\Z_8^2$.
\end{proposition}
By \cref{lem:pds}, the $(64,22,6)$ cases would follow from a regular Latin-square-type
$(64,21,8,6)$ partial difference set in the group, and the $(64,46,30)$ cases from the complement
of a regular negative-Latin-square-type $(64,18,2,6)$ partial difference set. An exhaustive search
over all symmetric subsets of the right size (a pseudo-Boolean formulation solved with z3,
\texttt{data/pds\_z3.py}; every nonexistence verdict was cross-checked by an independent
backtracking program, \texttt{data/pds\_ex.c} and \texttt{data/pds\_ex2.c}, with logs in
\texttt{data/}) shows that, among the eleven abelian groups
of order 64, regular $(64,21,8,6)$ partial difference sets and regular $(64,18,2,6)$ partial
difference sets exist in exactly the same three groups, $\Z_2^6$, $\Z_2^2\times\Z_4^2$ and
$\Z_8^2$, and in no other abelian group of order 64. (For $\Z_8^2$ both kinds appear, up to complementation, in the
enumeration of all regular partial difference sets in $C_8\times C_8$ by Malandro and
Smith~\cite[Table~5]{MalandroSmith2020}; for $\Z_2^6$ and $H\times H$ the Latin-square-type sets
are the partial congruence partitions of \cref{prop:known}.)
Consequently, the $(64,46,30)$ cases in $\Z_2^2\times\Z_4^2$ and in $\Z_8^2$ do follow from
\cref{lem:pds}, and we do not claim them as new, whereas the remaining nine cases of this
proposition cannot arise from \cref{lem:pds}; to our knowledge they are new in the same sense as
the cases of \cref{thm:new}. We keep them in a separate tier because the search for the partial
difference sets is a computation of a different kind from the rest of the paper. Our certificates
settle all eleven cases independently of this discussion.

\subsection{Database corrections}

\begin{proposition}\label{prop:known}
The following eight cases are listed as open in the repository, but follow from published
constructions or from their routine variants. We do not claim them as new; Lean-certified
examples are in \cref{tab:certs} (tier~c).
\begin{enumerate}
\item $(27,14,5)$ in $\Z_3^3$. Use \cite[Prop.~2.3]{Gordon2023}: $P$ is the Paley $(27,13,6)$
difference set of $\GF(27)$ and $N=\{0\}$. The group is $(\GF(27),+)\cong\Z_3^3$.
\item $(27,17,8)$ in $\Z_3^3$. Use \cite[Thm.~4.1]{Gordon2023} with $q=3$, $r=9$, so
$G=\GF(3)\oplus\GF(9)\cong\Z_3^3$. The repository attaches the comment ``Prime Power Pair'' to
the $\Z_3\times\Z_9$ entry. A $(27,17,8)$ SDS does also exist in $\Z_3\times\Z_9$; our
exhaustive search finds 1836 normalized ones.
\item $(27,10,1)$ in $\Z_3^3$. Take the product construction of \cite[Thm.~4.1]{HeChenGe2024}
with $m=1$, where it is stated only for even $m$. With $G=\Z_3\times\Z_3\times\Z_3$, take
$D'=\{1\}-\{2\}$, the $(3,2,-1)$ SDS of \cite[Prop.~2.2]{Gordon2023}, and set
$D=(1,0,\Z_3)+(0,\Z_3,0)+(0,D',D')$ (the parameter $x_1$ of \cite{HeChenGe2024} taken to be $0$).
The character computation of \cite{HeChenGe2024} applies verbatim, and we also checked it
directly.
\item $(36,11,2)$ in $\Z_6^2$. Apply \cref{lem:pds} to the rook-graph partial difference set
$(\Z_6\times0\cup 0\times\Z_6)\setminus\{0\}$ with parameters $(36,10,4,2)$. In $\Z_3\times\Z_{12}$
no $(36,10,4,2)$ partial difference set exists (exhaustive search), and that case is in
\cref{thm:new}.
\item $(64,22,6)$ in $\Z_2^6$. This is \cite[Thm.~5.3, Ex.~5.4]{HeWu2026}; it is not yet
reflected in the repository snapshot.
\item[6--7.] $(64,22,6)$ in $\Z_2^2\times\Z_4^2$ and in $\Z_8^2$. Apply \cref{lem:pds} to the
partial congruence partition partial difference set $(H\times0\cup0\times H\cup\{(h,h)\})\setminus\{0\}$
with parameters $(64,21,8,6)$ in $H\times H$, for $H=\Z_2\times\Z_4$ and $H=\Z_8$ respectively.
\item[8.] $(64,46,30)$ in $\Z_2^6$. Let $S$ be the set of nonzero vectors of $\GF(4)^3$ that
lie on the six points of a hyperoval of $\mathrm{PG}(2,4)$. Then $S$ is a $(64,18,2,6)$ partial
difference set. Its complement $G\setminus(S\cup\{0\})$ is a $(64,45,32,30)$ partial difference
set, and \cref{lem:pds} gives the SDS. We verified this computationally
(\texttt{data/hyperoval.py}). Our search independently found an example of exactly this
shape.
\end{enumerate}
\end{proposition}

%======================================================================
\section{Nonexistence for the remaining cases with \texorpdfstring{$v\le 28$}{v <= 28}}\label{sec:nonex}

\begin{proposition}[parity]\label{prop:parity}
Let $k=v-2$ with $v\ge 4$. If $v\not\equiv\lambda\pmod 2$, then no $(v,k,\lambda)$ SDS exists in
any abelian group of order $v$. In particular, no $(27,25,16)$ SDS exists in $\Z_{27}$ or in
$\Z_3\times\Z_9$.
\end{proposition}
\begin{proof}
Let $\{a,b\}$ be the zero set and take $y\notin\{0,a-b,b-a\}$. The pairs $(g,g+y)$ with both
coefficients nonzero are exactly those with $g\notin\{a,b,a-y,b-y\}$, and these four elements are
distinct. So $C_D(y)$ is a sum of $v-4$ terms $\pm1$, and $C_D(y)\equiv v\pmod 2$.
\end{proof}
He and Wu's small-defect criterion~\cite[Thm.~4.7]{HeWu2026} excludes $(27,25,16)$ only in
$\Z_3^3$.

\begin{proposition}[computer search, not formally verified]\label{prop:computer}
There is no SDS in any of the following 16 cases:
$(20,11,2)$ and $(20,17,8)$ in $\Z_2\times\Z_{10}$;
$(24,18,2)$ in $\Z_2\times\Z_{12}$ and in $\Z_2^2\times\Z_6$;
$(27,17,4)$, $(27,22,3)$ and $(27,22,9)$ in $\Z_3^3$ and in $\Z_3\times\Z_9$;
$(27,17,8)$ in $\Z_{27}$;
$(27,25,16)$ in $\Z_{27}$ and in $\Z_3\times\Z_9$;
$(28,13,4)$, $(28,19,10)$ and $(28,27,2)$ in $\Z_2\times\Z_{14}$.
\end{proposition}

The search is described in \cref{sec:methods}, and the log is \texttt{data/ex\_small.txt}. Each
case took between $0.01$ and $101$ seconds. \Cref{tab:small} summarizes all 23 open cases with
$v\le28$.

\begin{table}[ht]
\centering\small
\begin{tabular}{lll}
\toprule
case & status & source\\\midrule
$(9,8,1)$, $\Z_3^2$ & none & \cref{thm:981} (Lean)\\
$(20,11,2)$, $(20,17,8)$, $\Z_2\times\Z_{10}$ & none & \cref{prop:computer}\\
$(24,18,2)$, $\Z_2\times\Z_{12}$ and $\Z_2^2\times\Z_6$ & none & \cref{prop:computer}\\
$(25,12,1)$, $\Z_5^2$ & exists & \cref{prop:cyc25} (Lean)\\
$(27,10,1)$, $\Z_3\times\Z_9$; $(27,14,5)$, $\Z_3\times\Z_9$ & exist & \cref{thm:new} (Lean)\\
$(27,10,1)$, $(27,14,5)$, $(27,17,8)$, $\Z_3^3$ & exist & \cref{prop:known} (Lean)\\
$(27,17,4)$, $(27,22,3)$, $(27,22,9)$, $\Z_3^3$ and $\Z_3\times\Z_9$ & none & \cref{prop:computer}\\
$(27,17,8)$, $\Z_{27}$ & none & \cref{prop:computer}\\
$(27,25,16)$, $\Z_{27}$ and $\Z_3\times\Z_9$ & none & \cref{prop:parity}, \cref{prop:computer}\\
$(28,13,4)$, $(28,19,10)$, $(28,27,2)$, $\Z_2\times\Z_{14}$ & none & \cref{prop:computer}\\
\bottomrule
\end{tabular}
\caption{All 23 group-specific cases with $v\le 28$ that are open after \cite{HeWu2026}.}
\label{tab:small}
\end{table}

%======================================================================
\section{Search methods}\label{sec:methods}

\paragraph{Local search.} We store $s:G\to\{-1,0,1\}$ with $|P|=(k+s)/2$ and $|N|=(k-s)/2$
fixed, and minimize $\sum_{y\ne0}(C(y)-\lambda)^2$ by simulated annealing. Each move swaps two
coefficients. Changing $s_g$ by $d$ changes $C(y)$ by $d(s_{g+y}+s_{g-y})$, so an update costs
$O(v)$ (\texttt{data/sds\_ls.c}). All 710 open cases with $v\le 100$ were run for 4 CPU-seconds
each. Parameter triples that succeeded in at least one group of their order were retried for
20 seconds in the remaining groups.

\paragraph{Exhaustive search.} By \cref{lem:basic}(2) we may assume that $0\in S=\supp D$ and that
$s_0=+1$. We enumerate supports $S\ni 0$ of size $k$ and keep those satisfying the support parity
condition of \cref{lem:basic}(3). For each surviving support we then assign signs by
backtracking. For each shift $y$, the number of equal-sign pairs among the $c_y$ ordered pairs must
be $(c_y+\lambda)/2$ and the number of opposite-sign pairs $(c_y-\lambda)/2$. Partial counts are
pruned against these targets (\texttt{data/sds\_ex.c}). As a validation, the program finds
solutions in every case we tested where an SDS is known to exist: Gordon's $(18,13,4)$ in
$\Z_2\times\Z_3^2$, cyclic $(20,11,2)$, $(13,4,1)$ and $(7,6,-1)$, and our $(25,12,1)$,
$(27,10,1)$, $(27,14,5)$ and $(27,17,8)$. It reports none for the cases recorded as nonexistent,
namely $(18,13,4)$ in $\Z_{18}$ and $(9,8,1)$ in $\Z_9$ (\texttt{data/validate\_ex.txt}).

%======================================================================
\section{Formal verification}\label{sec:lean}

The file \texttt{lean/Main.lean} is self-contained and uses Lean~4.34.1 with core only, no
Mathlib. The group $\Z_{m_1}\times\dots\times\Z_{m_r}$ (moduli list \texttt{ms}, $v=\prod m_i$)
is modelled on $\{0,\dots,v-1\}$ by mixed-radix coordinates, most significant first. The group
law is componentwise addition modulo $m_i$:
\begin{lstlisting}
def coords : List Nat → Nat → List Nat
  | [], _ => []
  | m :: ms, x => (x / ms.prod) % m :: coords ms x
def ofCoords : List Nat → List Nat → Nat
  | [], _ => 0
  | _ :: ms, c => c.headD 0 * ms.prod + ofCoords ms c.tail
def addCoords : List Nat → List Nat → List Nat → List Nat
  | m :: ms, a :: as, b :: bs => (a + b) % m :: addCoords ms as bs
  | _, _, _ => []
def gadd (ms : List Nat) (x y : Nat) : Nat :=
  ofCoords ms (addCoords ms (coords ms x) (coords ms y))

def IsSDS (ms : List Nat) (k : Nat) (lam : Int) (D : Nat → Int) : Prop :=
  (∀ g ∈ List.range ms.prod, D g = -1 ∨ D g = 0 ∨ D g = 1) ∧
  ((List.range ms.prod).filter (fun g => D g ≠ 0)).length = k ∧
  ∀ y ∈ List.range ms.prod, y ≠ 0 →
    ((List.range ms.prod).map (fun g => D g * D (gadd ms g y))).sum = lam
\end{lstlisting}
The main statements are as follows; there is one existence theorem for each of the 45 cases:
\begin{lstlisting}
theorem no_SDS_9_8_1_Z3xZ3 : ¬ ∃ D : Nat → Int, IsSDS [3,3] 8 1 D
theorem sds_25_12_1_G5x5   : ∃ D : Nat → Int, IsSDS [5, 5] 12 (1) D
theorem sds_27_10_1_G3x9   : ∃ D : Nat → Int, IsSDS [3, 9] 10 (1) D
...                        -- 45 theorems sds_v_k_lambda_G..., see Appendix A
theorem model_G5x5 : GroupModelOK [5, 5]   -- one sanity lemma per group (23 groups)
\end{lstlisting}

\paragraph{Statement audit.}
\texttt{GroupModelOK ms} states two things. First, \texttt{coords} maps $\{0,\dots,v-1\}$ into
valid tuples and has left inverse \texttt{ofCoords}. Second, \texttt{ofCoords} maps every tuple
of the Cartesian product $[0,m_1)\times\dots\times[0,m_r)$ into $\{0,\dots,v-1\}$ and has left
inverse \texttt{coords}. So \texttt{coords} is a bijection. Since \texttt{gadd} is defined as
transported componentwise addition, $(\{0,\dots,v-1\},\texttt{gadd ms})$ is exactly the group
with invariant factors \texttt{ms}, which is the repository label. \texttt{IsSDS} says three
things: the coefficients lie in $\{-1,0,1\}$; exactly $k$ of them are nonzero; and
\eqref{eq:autocorr} holds for every nonzero $y\in G$, with every sum running once over all of
$G$. This is exactly \eqref{eq:sds}. The existence theorems are proved by exhibiting $D$, built
from the explicit lists $P,N$ of coordinate tuples, followed by \texttt{decide +kernel}. Whether
the helper \texttt{signedOf} that builds $D$ is correct does not matter, because the predicate
is checked on the resulting function itself.

The nonexistence theorem quantifies over \emph{all} $D:\mathbb N\to\Z$. Its proof has three
steps. First, the three blocks $[D\,0,D\,1,D\,2]$, $[D\,3,D\,4,D\,5]$ and $[D\,6,D\,7,D\,8]$ lie in
the list of all 27 ternary words of length~3. Second, \texttt{IsSDS [3,3]} depends only on these
nine values. Third, 27 kernel-checked lemmas exclude all $3^9=19683$ coefficient vectors. No
translation or negation normalization is used.

\paragraph{Axioms and resources.} \texttt{lean/lean\_output.txt} contains 69 \texttt{\#print axioms}
lines. The output is:
\begin{lstlisting}
'no_SDS_9_8_1_Z3xZ3' depends on axioms: [propext, Quot.sound]
'model_G5x5' depends on axioms: [propext]        -- same for all 23 model lemmas
'sds_25_12_1_G5x5' depends on axioms: [propext, Quot.sound]
                                                 -- same for all 45 existence theorems
leanchk: wall=1221.58s maxrss=5790648kB exit=0
\end{lstlisting}
There is no \texttt{sorryAx}, no \texttt{native\_decide} (no \texttt{Lean.ofReduceBool}), and
there are no \texttt{axiom} declarations, errors or warnings. The last line is the resource log
of the wrapper that ran Lean. To reproduce, run \texttt{lean Main.lean} in the directory
\texttt{lean/} with a Lean~4.34.1 toolchain; no project file and no Mathlib are needed. The
recorded run took 1222 seconds of wall time and 5.8\,GB of memory; an independent re-run during
review (same file, SHA-256
\texttt{5a0ae726\dots3ce4e27}) took 587 seconds and 5.8\,GB and produced identical output.

%======================================================================
\section{Limitations and open questions}\label{sec:limits}

\begin{itemize}
\item \emph{Novelty.} ``New'' means: open in the repository (v1.2) and after \cite{HeWu2026},
and not implied by any construction we could identify in
\cite{Gordon2023,HeChenGe2024,HeWu2026} or by \cref{lem:pds}. For \cref{lem:pds} the question is
settled by exhaustive search in all abelian groups of order 64 (\cref{prop:maybe}); the only
other certificates with $|N|=1$ are $(27,14,5)$ in $\Z_3\times\Z_9$ and $(36,11,2)$ in
$\Z_3\times\Z_{12}$, and there \cref{lem:pds} cannot apply either: a $(27,13,7,5)$ partial
difference set would be a strongly regular graph with non-integral eigenvalue multiplicities
($6f=39$), and no $(36,10,4,2)$ partial difference set exists in $\Z_3\times\Z_{12}$
(\cref{prop:known}). We did not survey the partial difference set or Hadamard difference set
literature exhaustively, so some 2-group examples might still follow from constructions there.
Finally, the group-specific open cases themselves date only from version~1.2 of the repository
(April 2026); earlier versions covered essentially only cyclic groups.
\item \cref{prop:computer} is the output of a C program, validated as described in
\cref{sec:methods}. It is not formally verified. Only \cref{thm:981} and the 45 existence
results are Lean-certified.
\item Local search found nothing for the 455 open cases with $65\le v\le100$. Nothing should be
inferred from this.
\item Open questions. Is there a uniform construction behind the $n=16$ families of
\cref{thm:new}? Do $(32,20,4)$ and $(32,28,12)$ exist in the remaining groups of order~32
(including $\Z_{32}$)? Do $(48,24,8)$ and $(48,32,16)$ exist in $\Z_2\times\Z_{24}$ and $\Z_{48}$?
Are there further eighth-order cyclotomic SDSs in $\GF(p^2)$?
\end{itemize}

%======================================================================
\begin{thebibliography}{9}
\bibitem{Gordon2023}
D.~M.~Gordon, \emph{Signed difference sets}, Des.\ Codes Cryptogr.\ \textbf{91} (2023), no.~5, 2107--2115;
\href{https://arxiv.org/abs/2212.10630}{arXiv:2212.10630}.
\bibitem{GordonSDSrepo}
D.~M.~Gordon, \emph{La Jolla Signed Difference Set Repository}, version 1.2 (April 24, 2026),
\texttt{sds.json}, Zenodo, concept \href{https://doi.org/10.5281/zenodo.7473882}{doi:10.5281/zenodo.7473882},
version \href{https://doi.org/10.5281/zenodo.19741004}{doi:10.5281/zenodo.19741004};
source \url{https://github.com/dmgordo/signed-difference-sets}; see also
\url{https://dmgordon.org/signed-difference-sets/}. Versions 1.0 (December 20, 2022) and 1.1
(March 9, 2023) are \href{https://doi.org/10.5281/zenodo.7473883}{doi:10.5281/zenodo.7473883} and
\href{https://doi.org/10.5281/zenodo.7713113}{doi:10.5281/zenodo.7713113}.
\bibitem{HeChenGe2024}
Z.~He, T.~Chen, G.~Ge, \emph{New constructions of signed difference sets}, Des.\ Codes Cryptogr.\
\textbf{92} (2024), no.~8, 2323--2340; \href{https://arxiv.org/abs/2306.05631}{arXiv:2306.05631}.
\bibitem{HeWu2026}
Z.~He, J.~Wu, \emph{Nonexistence results and constructions for signed difference sets},
\href{https://arxiv.org/abs/2608.22936}{arXiv:2608.22936v2} (September 20, 2026).
\bibitem{HeScript}
Z.~He, \emph{SDS Open-Case Verification} (Python code and data),
\url{https://github.com/zhiwenhe94/sds-open-case-verification}, accessed October 8, 2026.
\bibitem{MalandroSmith2020}
M.~E.~Malandro, K.~W.~Smith, \emph{Partial difference sets in $C_{2^n}\times C_{2^n}$}, Discrete
Math.\ \textbf{343} (2020), no.~4, 111744; \href{https://arxiv.org/abs/1811.11223}{arXiv:1811.11223}.
\end{thebibliography}

%======================================================================
\appendix
\section{All certificates}\label{app:certs}

Elements of $\Z_{m_1}\times\dots\times\Z_{m_r}$ are written as mixed-radix indices: the tuple
$(x_1,\dots,x_r)$ with $0\le x_i<m_i$ is the integer
$x=x_1 m_2\cdots m_r+x_2m_3\cdots m_r+\dots+x_r$. For example, in $\Z_3\times\Z_9$ the index $16$
is $(1,7)$. The same tuples, together with the indices, are in \texttt{data/certificates.json},
and the tuples appear verbatim in \texttt{lean/Main.lean}. Tier: a = new (\cref{thm:new});
b = order-64 cases of \cref{prop:maybe} (nine of them cannot come from \cref{lem:pds}, two do);
c = database correction (\cref{prop:known}).

{\small
\setlength{\tabcolsep}{4pt}
\begin{longtable}{@{}l l c p{7.2cm} p{2.6cm}@{}}
\caption{Explicit signed difference sets $D=P-N$ (indices of $P$ and of $N$).}\label{tab:certs}\\
\toprule
$(v,k,\lambda)$ & group & tier & $P$ & $N$\\\midrule
\endfirsthead
\toprule
$(v,k,\lambda)$ & group & tier & $P$ & $N$\\\midrule
\endhead
\bottomrule
\endfoot
\certrow{(25,12,1)}{$\Z_{5}^{2}$}{a}{1, 5, 11, 12, 14, 16, 17, 19, 20}{10, 22, 23}
\certrow{(27,10,1)}{$\Z_{3}^{3}$}{c}{1, 11, 13, 16, 17, 22, 24, 25}{5, 10}
\certrow{(27,10,1)}{$\Z_{3}\times\Z_{9}$}{a}{6, 7, 12, 16, 18, 21, 24, 25}{3, 15}
\certrow{(27,14,5)}{$\Z_{3}^{3}$}{c}{0, 1, 3, 9, 11, 13, 14, 16, 18, 21, 22, 25, 26}{8}
\certrow{(27,14,5)}{$\Z_{3}\times\Z_{9}$}{a}{0, 1, 3, 12, 14, 15, 16, 17, 18, 22, 23, 24, 26}{2}
\certrow{(27,17,8)}{$\Z_{3}^{3}$}{c}{0, 1, 2, 3, 4, 5, 6, 12, 13, 15, 17, 19, 22, 23, 24, 25}{26}
\certrow{(32,20,4)}{$\Z_{2}^{5}$}{a}{2, 4, 6, 7, 8, 10, 11, 12, 13, 16, 17, 19, 20, 22, 25, 29}{15, 23, 28, 31}
\certrow{(32,20,4)}{$\Z_{2}^{3}\times\Z_{4}$}{a}{2, 3, 6, 9, 10, 13, 15, 16, 18, 19, 20, 21, 22, 24, 26, 27}{1, 11, 12, 31}
\certrow{(32,20,4)}{$\Z_{2}^{2}\times\Z_{8}$}{a}{3, 8, 9, 11, 13, 15, 18, 19, 20, 22, 23, 25, 26, 27, 29, 30}{7, 12, 16, 31}
\certrow{(32,20,4)}{$\Z_{2}\times\Z_{4}^{2}$}{a}{0, 1, 2, 3, 6, 8, 9, 10, 13, 15, 16, 19, 20, 21, 24, 30}{4, 11, 27, 31}
\certrow{(32,20,4)}{$\Z_{2}\times\Z_{16}$}{a}{0, 1, 3, 4, 5, 6, 8, 11, 12, 14, 18, 19, 21, 25, 26, 27}{9, 13, 17, 29}
\certrow{(32,20,4)}{$\Z_{4}\times\Z_{8}$}{a}{0, 1, 2, 6, 7, 8, 10, 13, 16, 17, 19, 20, 24, 25, 29, 30}{4, 9, 11, 15}
\certrow{(32,28,12)}{$\Z_{2}\times\Z_{4}^{2}$}{a}{2, 3, 4, 5, 6, 8, 9, 10, 11, 13, 14, 15, 17, 19, 20, 21, 22, 23, 25, 27, 28, 29, 30, 31}{0, 1, 7, 12}
\certrow{(32,28,12)}{$\Z_{4}\times\Z_{8}$}{a}{1, 2, 4, 5, 6, 8, 9, 10, 11, 12, 13, 14, 15, 16, 17, 18, 20, 21, 25, 27, 28, 29, 30, 31}{0, 22, 24, 26}
\certrow{(36,11,2)}{$\Z_{3}\times\Z_{12}$}{a}{2, 5, 9, 11, 13, 17, 18, 23, 34, 35}{29}
\certrow{(36,11,2)}{$\Z_{6}^{2}$}{c}{5, 6, 7, 9, 13, 22, 27, 31, 32, 33}{23}
\certrow{(48,24,8)}{$\Z_{2}^{3}\times\Z_{6}$}{a}{0, 1, 3, 5, 10, 11, 12, 14, 19, 20, 21, 22, 25, 30, 31, 36, 37, 39, 42, 43, 45, 46}{24, 40}
\certrow{(48,24,8)}{$\Z_{2}^{2}\times\Z_{12}$}{a}{2, 3, 4, 6, 9, 11, 12, 13, 15, 16, 19, 20, 21, 23, 28, 29, 30, 32, 36, 38, 40, 41}{1, 7}
\certrow{(48,24,8)}{$\Z_{4}\times\Z_{12}$}{a}{1, 2, 4, 5, 7, 16, 17, 19, 23, 24, 26, 27, 28, 30, 33, 34, 35, 37, 40, 41, 44, 47}{10, 38}
\certrow{(48,32,16)}{$\Z_{2}^{3}\times\Z_{6}$}{a}{0, 1, 2, 4, 5, 6, 7, 10, 15, 18, 19, 20, 23, 24, 26, 29, 30, 31, 32, 34, 35, 36, 38, 40, 41, 43, 44, 45, 46, 47}{22, 37}
\certrow{(48,32,16)}{$\Z_{2}^{2}\times\Z_{12}$}{a}{0, 1, 3, 5, 6, 10, 11, 12, 15, 18, 19, 21, 22, 24, 26, 27, 29, 30, 31, 32, 33, 34, 35, 36, 37, 38, 42, 44, 45, 46}{9, 39}
\certrow{(48,32,16)}{$\Z_{4}\times\Z_{12}$}{a}{1, 3, 4, 5, 7, 8, 9, 10, 11, 12, 13, 14, 16, 19, 21, 22, 23, 24, 29, 30, 33, 35, 37, 38, 40, 41, 42, 43, 45, 46}{2, 27}
\certrow{(54,46,10)}{$\Z_{3}^{2}\times\Z_{6}$}{a}{0, 1, 3, 4, 6, 9, 12, 13, 14, 16, 21, 22, 23, 25, 26, 29, 30, 32, 33, 34, 35, 38, 40, 41, 42, 44, 45, 46, 47, 48, 49, 50, 51, 52, 53}{7, 10, 15, 17, 18, 19, 20, 28, 31, 37, 43}
\certrow{(64,18,2)}{$\Z_{2}^{6}$}{a}{0, 5, 7, 8, 11, 24, 25, 27, 31, 32, 33, 37, 41, 47, 52}{4, 39, 49}
\certrow{(64,18,2)}{$\Z_{2}^{4}\times\Z_{4}$}{a}{2, 6, 8, 10, 18, 23, 24, 27, 32, 41, 42, 43, 45, 54, 62}{35, 52, 58}
\certrow{(64,18,2)}{$\Z_{2}^{3}\times\Z_{8}$}{a}{2, 4, 17, 18, 21, 24, 25, 26, 28, 29, 34, 36, 48, 50, 62}{20, 30, 58}
\certrow{(64,18,2)}{$\Z_{2}^{2}\times\Z_{4}^{2}$}{a}{0, 1, 5, 7, 8, 10, 12, 17, 18, 34, 43, 44, 48, 52, 62}{14, 38, 57}
\certrow{(64,18,2)}{$\Z_{2}\times\Z_{4}\times\Z_{8}$}{a}{5, 8, 13, 16, 17, 21, 22, 23, 26, 37, 43, 44, 50, 52, 62}{15, 41, 51}
\certrow{(64,18,2)}{$\Z_{4}^{3}$}{a}{3, 7, 16, 20, 21, 23, 25, 26, 28, 30, 35, 36, 44, 50, 52}{18, 19, 47}
\certrow{(64,18,2)}{$\Z_{8}^{2}$}{a}{9, 14, 15, 16, 20, 25, 27, 28, 31, 41, 42, 43, 50, 54, 56}{13, 45, 57}
\certrow{(64,22,6)}{$\Z_{2}^{6}$}{c}{2, 6, 7, 8, 10, 11, 13, 14, 18, 20, 22, 25, 37, 38, 39, 45, 48, 49, 54, 56, 63}{57}
\certrow{(64,22,6)}{$\Z_{2}^{4}\times\Z_{4}$}{b}{3, 7, 8, 10, 11, 12, 14, 15, 18, 23, 25, 28, 32, 35, 41, 42, 49, 51, 59, 60, 61}{56}
\certrow{(64,22,6)}{$\Z_{2}^{3}\times\Z_{8}$}{b}{0, 2, 6, 7, 8, 11, 14, 16, 21, 24, 26, 27, 31, 32, 34, 35, 36, 41, 47, 53, 54}{13}
\certrow{(64,22,6)}{$\Z_{2}^{2}\times\Z_{4}^{2}$}{c}{5, 6, 19, 21, 27, 28, 29, 31, 35, 38, 41, 44, 45, 46, 49, 50, 54, 58, 59, 61, 63}{53}
\certrow{(64,22,6)}{$\Z_{2}\times\Z_{4}\times\Z_{8}$}{b}{0, 5, 8, 9, 12, 15, 17, 18, 20, 22, 25, 31, 32, 33, 34, 35, 38, 40, 44, 51, 52}{37}
\certrow{(64,22,6)}{$\Z_{4}^{3}$}{b}{2, 7, 8, 10, 18, 26, 27, 30, 32, 33, 34, 37, 40, 43, 44, 46, 49, 58, 60, 61, 63}{56}
\certrow{(64,22,6)}{$\Z_{4}\times\Z_{16}$}{b}{1, 2, 4, 6, 7, 10, 11, 14, 19, 20, 25, 31, 33, 35, 36, 47, 52, 53, 55, 59, 61}{49}
\certrow{(64,22,6)}{$\Z_{8}^{2}$}{c}{1, 2, 3, 4, 5, 6, 7, 8, 9, 16, 18, 24, 27, 32, 36, 40, 45, 48, 54, 56, 63}{0}
\certrow{(64,46,30)}{$\Z_{2}^{6}$}{c}{0, 1, 2, 3, 5, 7, 8, 9, 11, 12, 14, 15, 16, 17, 18, 20, 22, 23, 24, 25, 26, 27, 28, 31, 32, 33, 34, 35, 36, 37, 38, 39, 40, 44, 46, 47, 48, 50, 51, 53, 56, 57, 58, 60, 61}{59}
\certrow{(64,46,30)}{$\Z_{2}^{4}\times\Z_{4}$}{b}{0, 1, 2, 3, 4, 5, 6, 7, 8, 9, 11, 12, 14, 15, 17, 18, 20, 21, 22, 23, 24, 25, 26, 31, 32, 33, 37, 38, 41, 42, 43, 44, 46, 47, 48, 49, 50, 51, 52, 53, 56, 57, 58, 60, 62}{63}
\certrow{(64,46,30)}{$\Z_{2}^{3}\times\Z_{8}$}{b}{1, 2, 3, 6, 7, 8, 9, 10, 11, 15, 16, 19, 21, 22, 23, 24, 26, 27, 28, 29, 31, 32, 33, 35, 36, 37, 38, 39, 40, 41, 42, 43, 45, 46, 47, 48, 50, 54, 56, 57, 59, 60, 61, 62, 63}{30}
\certrow{(64,46,30)}{$\Z_{2}^{2}\times\Z_{4}^{2}$}{b}{0, 1, 3, 4, 5, 7, 9, 10, 11, 12, 14, 15, 16, 17, 18, 20, 21, 23, 24, 26, 27, 32, 34, 35, 36, 37, 38, 40, 42, 43, 44, 45, 47, 48, 49, 50, 52, 53, 55, 56, 57, 58, 60, 61, 62}{29}
\certrow{(64,46,30)}{$\Z_{2}\times\Z_{4}\times\Z_{8}$}{b}{4, 6, 7, 8, 9, 10, 11, 12, 14, 15, 16, 17, 21, 22, 23, 24, 26, 28, 29, 30, 32, 33, 34, 35, 36, 37, 39, 40, 41, 43, 44, 45, 46, 47, 48, 49, 50, 52, 53, 55, 56, 58, 59, 60, 63}{51}
\certrow{(64,46,30)}{$\Z_{4}^{3}$}{b}{0, 1, 2, 3, 5, 6, 9, 10, 11, 13, 14, 15, 16, 17, 18, 19, 20, 21, 22, 23, 24, 27, 29, 30, 31, 32, 33, 34, 35, 36, 38, 41, 42, 43, 44, 46, 47, 48, 49, 53, 55, 57, 58, 59, 62}{26}
\certrow{(64,46,30)}{$\Z_{8}^{2}$}{b}{0, 1, 2, 3, 4, 5, 6, 7, 8, 9, 10, 11, 13, 14, 18, 21, 22, 23, 25, 26, 27, 28, 30, 31, 32, 34, 36, 38, 39, 40, 41, 42, 45, 46, 47, 48, 49, 52, 55, 58, 59, 60, 61, 62, 63}{35}

\end{longtable}
}

\end{document}
